% Generated draft: requires mathematical and source review.
\documentclass[11pt]{article}
\usepackage[margin=25mm]{geometry}
\usepackage{fontspec}
\setmainfont{Latin Modern Roman}
\usepackage{amsmath,amssymb,mathtools,unicode-math}
\setmathfont{Latin Modern Math}
\usepackage{enumitem,needspace,xurl,hyperref}
\setlength{\parindent}{0pt}
\setlength{\parskip}{5pt}
\setlength{\emergencystretch}{3em}
\newcommand{\N}{\mathbb{N}}
\newcommand{\Z}{\mathbb{Z}}
\newcommand{\Q}{\mathbb{Q}}
\newcommand{\R}{\mathbb{R}}
\newcommand{\C}{\mathbb{C}}
\newcommand{\F}{\mathbb{F}}
\newcommand{\A}{\mathbb{A}}
\newcommand{\PP}{\mathbb{P}}
\newcommand{\E}{\mathbb{E}}
\newcommand{\eps}{\varepsilon}
\newcommand{\dd}{\,\mathrm{d}}
\begin{document}
\section*{UnsolvedMath ID: 30002606}
\textbf{Source problem code:} catalogue entry
\subsection{Definitions and mathematical statement}
Let \(\delta>0\) denote the peridynamic interaction horizon, and let \(h_j>0\) denote the local mesh size associated with element \(j\). These are the interpretations of the notation in the supplied record. Let \(J_{\mathrm{near}}\) and \(J_{\mathrm{away}}\) denote the element indices classified as near defects and away from defects, respectively; the record does not specify these classifications or whether they exhaust the mesh.

The recorded construction request is to develop an adaptive Galerkin discretization of a single peridynamic model, used as a computational multiscale mono-model, satisfying
\[
\begin{aligned}
&h_j<\delta &&\text{with discontinuous Galerkin elements for every }j\in J_{\mathrm{near}},\\
&h_j>\delta &&\text{with continuous Galerkin elements for every }j\in J_{\mathrm{away}},
\end{aligned}
\]
together with defect detection, adaptive refinement, and suitable numerical quadrature.

This specifies a proposed construction rather than a complete mathematical existence or convergence problem. The peridynamic equation, spatial dimension and domain, interaction kernel or force law, data, boundary or volume constraints, definition of a defect, classification and refinement criteria, Galerkin spaces and their coupling, quadrature requirements, and desired accuracy or stability conclusions are not supplied. The case \(h_j=\delta\) and treatment of transition elements are also unspecified. Consequently, no fully quantified correctness or convergence assertion can be recovered from this record alone.

The supplied statement-verification note says that the report presents these finite-element choices as a proposed computational construction and reports one-dimensional results, rather than posing this extraction as a standalone open problem. The registry marks the statement as \texttt{unrecoverable} and supplies no clean replacement statement. Its dated assessment describes higher-dimensional implementation and analysis as incomplete; this is supplied curation, not an independently verified mathematical conclusion.
\subsection{Short English statement}
The recorded aim is an adaptive discretization of one peridynamic model: use discontinuous Galerkin elements with mesh size below the interaction horizon near defects, and continuous Galerkin elements with mesh size above the horizon away from defects, together with detection, refinement, and quadrature. The supplied verification describes this as a proposed construction with one-dimensional results. Missing model specifications and success criteria prevent recovery of a precise standalone open problem.
\subsection{Sources}
\begin{itemize}
\item Local UnsolvedMath registry, permanent ID \(30002606\), title \emph{Adaptive Galerkin Discretization of Peridynamics}, exact source problem code \texttt{catalogue entry}, in the Oberwolfach Reports Open Problems collection. The record gives proposed year 2014 and names no proposer.
\item Supplied source citation: \emph{Computational Multiscale Methods (2015)}, Oberwolfach Reports, DOI 10.4171/owr/2014/30, pp.\ 1625--1681. \url{https://doi.org/10.4171/owr/2014/30}. The background attributes the material to workshop and open-problem discussions documented in Oberwolfach Reports.
\item Supplied literature source: \emph{Oberwolfach Report 30/2014}. \url{https://ems.press/content/serial-article-files/46520}.
\item Literature assessment date: \texttt{2026-08-21}. This is supplied dated curation, not independently verified here. Its status label is \texttt{partially\_solved}; its assessment describes a concrete mono-model recipe involving horizon-relative mesh size, adaptive refinement, continuous versus discontinuous finite elements, and one-dimensional analysis, with higher-dimensional implementation and analysis remaining incomplete. The background explicitly requests independent verification of time-sensitive status and cited claims.
\item The supplied statement-verification note reports checking against the official report and classifies the extraction as a proposed construction rather than a standalone open problem. The registry retains the original record for provenance, labels its statement \texttt{unrecoverable}, and promotes no clean replacement. These verification claims are reproduced as supplied metadata; no source visit or independent verification is claimed here.
\end{itemize}
\end{document}
